\chapter{波动光学}
\label{chap:波动光学}
光是我们最熟悉的自然现象。光学（optics）和力学、天文学、几何学一样，是在物理学中发展最早的几门学科之一。然而在相当长的时间中，我们对光学的知识仅局限在一些简单的成像规律上。随着天文学和生物学的发展，伴随光学仪器的研究和制造，逐步形成了以光线为基础，通过几何学的方法研究光在透明介质中传播规律的\uwave{几何光学}。几何光学最为我们熟知的重大成就莫过于望远镜和显微镜的发明，前者使得天文学得以飞跃发展，并间接导致了整个经典力学体系的创建，后者则直接导致了生物学中对细胞的发现。但几何光学只不过是对光现象的规律的一种客观描述，其并没有触及“光是什么”的本质问题。光在几何光学中常常以所谓“光线”的形态存在，但这只是一种数学上的建模方式罢了，其并未揭示光的存在形态。

光到底是什么？在17世纪的物理史上曾经存在过两种对立的观点
\begin{itemize}
    \item \uwave{光的微粒说}，以牛顿为代表的学者认为光是微粒。
    \item \uwave{光的波动说}，以惠更斯为代表的学者认为光是一种称为以太的特殊介质中的机械波动。
\end{itemize}
光的微粒说在最初占据了主流地位，但是由于光的干涉现象逐渐被发现，光的机械波动说逐渐得到了确立。而之后，麦克斯韦通过麦克斯韦方程组预言了光是一种电磁波，并且在之后由赫兹通过实验证实。而更进一步的，迈克耳孙-莫雷实验的零结果最终否定了作为电磁波介质的以太的存在性，证明了电磁波的传播不需要介质的，至此光的电磁理论被建立。尽管光在电磁理论的观点下仍是波动，但在实质上与最初的波动说已经完全不同了，称之为\uwave{波动光学}。

\input{chapter/Chapter13A.tex}
\section{光的干涉}
\label{sec:光的干涉}
在\xref{sec:光的电磁理论}中，我们在波动光学重新认识了光的基本形态，并引入了光强和光程这两个非常重要的概念。而本节我们就来思考这样一个问题，如果两束光在某处叠加将会导致什么？类比机械波的知识，我们知道，两列同频率且具有固定相位差的波将会发生干涉现象，而这种干涉现象是一切波动所具有的共同特性，无所谓是机械波还是电磁波，因此光也会发生干涉现象。

\subsection{光的相干性}
\label{subsec:光的相干性}
光并不总是会发生干涉现象的，而是要满足以下的\uwave{相干条件}
\begin{enumerate}
    \item 光的频率相同
    \item 光存在平行的振动分量
    \item 光在叠加时的相位差恒定
\end{enumerate}
我们将满足相干条件的光波称为\uwave{相干光波}，在上述三个条件中，第二条是比较费解的，这是因为我们知道，取决于光的传播方向，光矢量的振动方向是不一样的，试想若两束光的光矢量分别在$x$轴和$y$轴方向振动，即便它们频率相同且相位差恒定，那显然也无法产生干涉现象。

我们现在来思考，光的干涉将导致什么现象？我们可能觉得不会有什么现象，因为干涉不过是简谐振动的叠加，而简谐振动的叠加仍然是简谐振动。但问题在于，这种叠加导致了振动在空间上的分布，也就是波动的形态发生了变化，有些点的振动加强，有些点的振动减弱，而我们知道我们能看到的光的明暗实质是光强，而光强又与光振幅的平方成正比，因此我们就可以预期，光发生干涉时，有些地方会更亮，有些地方会更暗，\empx{光强在干涉时并不是简单的叠加}。

在这一小节，我们姑且先考虑光的干涉在某一点的光强，假设空间中有两束同频率、相位差恒定、光振动方向一致的光波，传播到空间中某点$P$处，记其分别引发的光振动为$E_A,E_B$，则
\gdef\prefixeq{光的相干性}
\begin{Gather}[6pt]
    E_\text{A}=E_{\text{A}0}\cos(\omega t+\phi_{\text{A}})\xlabelpeq{1}\\
    E_\text{B}=E_{\text{B}0}\cos(\omega t+\phi_{\text{A}})\xlabelpeq{2}
\end{Gather}

根据光矢量的叠加原理，在任意时刻都有
\begin{Equation}&[3]
    E=E_1+E_2
\end{Equation}
其中$E$为点$P$处光的合振动，而两个简谐振动的合成仍然为简谐振动
\begin{Equation}&[4]
    E=E_0\cos(\omega t+\phi)
\end{Equation}
其中上式光的合振动的振幅为\footnote{这里合振动的初相位$\phi$也是可以求出的，其为$\mal{\phi=\arctan\frac{E_{\text{A}0}\sin\phi_{\text{A}}-E_{\text{B}0}\sin\phi_{\text{B}}}{E_{\text{A}0}\cos\phi_{\text{A}}-E_{\text{B}0}\cos\phi_{\text{B}}}}$，但此处用不到该结果。}
\begin{Equation}&[5]
    E_0=\sqrt{E_{\text{A}0}^2+E_{\text{B}0}^2+2E_{\text{A}0}E_{\text{B}0}\cos(\phi_{\text{B}}-\phi_{\text{A}})}
\end{Equation}
而在这种情况下，根据\fancyref{fml:光的强度}
\begin{Equation}&[6]
    I=\frac{n}{\mu_0c}\frac{1}{\tau}\Int[0][\tau]{E^2\dd{t}}=\frac{n}{\mu_0c}\frac{1}{\tau}\Int[0][\tau]{E_0^2\cos(\omega t+\phi)\dd{t}}
\end{Equation}
这就得到，其中$\delt{\phi}=\phi_\text{B}-\phi_\text{A}$为相位差
\begin{Equation}&[7]
    I=\frac{n}{2\mu_0c}\qty[E_{\text{A}0}^2+E_{\text{B}0}^2+2E_{\text{A}0}E_{\text{B}0}\cos\delt{\phi}]
\end{Equation}
我们考虑到$E_{\text{A}0}^2(n/2\mu_0 c)=I_{\text{A}}$以及$E_{\text{B}0}^2(n/2\mu_0 c)=I_{\text{B}}$
\begin{OldBoxFormula}[简谐光的干涉光强]
    简谐相干光在发生干涉时的光强为，其中$\delt{\phi}$为相位差
    \begin{Equation}&[8]
        I=I_{\text{A}}+I_{\text{B}}+2\sqrt{I_\text{A}}\sqrt{I_\text{B}}\cos{\delt{\phi}}
    \end{Equation}
    这就是说，干涉时的光强是两束光的光强之和，再加上$2\sqrt{I_\text{A}}\sqrt{I_\text{B}}\cos{\delt{\phi}}$，称为\uwave{干涉项}。

    如果定义\uwave{光程差} $\delta=\Delta_\text{B}-\Delta_{A}$，那么考虑到$\delt{\phi}=k\delta=(2\pi/\lambda)\delta$
    \begin{Equation}
        I=I_{\text{A}}+I_{\text{B}}+2\sqrt{I_\text{A}}\sqrt{I_\text{B}}\cos(\frac{2\pi}{\lambda}\delta)
    \end{Equation}
    这样就将光强的干涉项与相位差的关系转化为与光程差的关系。
\end{OldBoxFormula}

由\xref{fml:简谐光的干涉光强}中我们不难看出相位差和光程差与\uwave{干涉加强处}和\uwave{干涉减弱处}的关系
\begin{OldTable}[干涉光强与相位差和光程差的关系]
<状态&\mc{2}{相位差}&\mc{2}{光程差}&取值\\>
干涉加强&$\pi$的偶数倍&$\delt{\phi}=\pm \pi(2m)$&半波长的偶数倍&$\delta=\pm(\lambda/2)(2m)$&$m\in\N$\\
干涉减弱&$\pi$的奇数倍&$\delt{\phi}=\pm \pi(2m-1)$&半波长的奇数倍&$\delta=\pm(\lambda/2)(2m-1)$&$m\in\N^{*}$\\
\end{OldTable}

特别的，如果参与干涉的两束光的光强相同，即$I_\text{A}=I_\text{B}=I_0$，那么\xref{fml:简谐光的干涉光强}有
\begin{Equation}
    I=2I_0+2I_0\cos(\frac{2\pi}{\lambda}\delta)
\end{Equation}
这时干涉加强处将有$I=4I_0$即达到单个光波光强的四倍，而干涉减弱处则有$I=0$。

\subsection{光源的发光机制}
\label{subsec:光源的发光机制}
在\xref{subsec:光的相干性}中我们讨论了光的干涉现象，但这总让人感觉有些奇怪，毕竟当我们打开桌上的台灯时，我们的桌子将被均匀的照亮，我们从未在该过程中看到过任何干涉现象，尽管从理论上来说灯管上两个不同位置的光源似乎是可以发生干涉的，这就与光源的发光机制有关了。

在化学中，我们曾经提到过，电子由激发态跃迁至基态时将以电磁波的形式放出能量，这也就是物质燃烧时火焰的光亮来源。而实际上，现代使用的各类照明设施，无论是节能灯（荧光灯和电子镇流器）还是LED灯（发光二极管），尽管用于激发电子的方式不同，但是就其发光的实质而言仍然是电子跃迁释放的可见光频段的电磁波\footnote{总之通过RLC电路的电磁振荡是无法产生可见光频段的电磁波的，因为可见光的频率太高了。}。而我们知道，尽管电子跃迁是大量且持续的，使得宏观上能产生稳定的发光现象，但是单个电子跃迁的持续发光时间都是非常短的，因此光波的初相位是不断变化的（每换一个电子就换一个初相位），而位于光源两个不同位置的光波往往是由两个不同的原子发射的，由于初相位均在不断变化，因此两者的光波在空间相遇时就不具有稳定的相位差，这样就无法形成干涉现象。在数学上就是$\cos{\delt{\phi}}$并不是一个无关时间的常数，不能从\xrefpeq[光的相干性]{6}中作为常数提出，而由于相位差$\delt{\phi}$的变化是完全随机的，因此$\cos{\delt{\phi}}$通过积分求出的均值为零，这样干涉项就为零了，此时的总光强即为
\begin{Equation}
    I=I_1+I_2
\end{Equation}
即光在非相干叠加时的光强就是两者光强的和，这个结果也很好理解，原先在干涉叠加时振动增强和振动减弱是稳定的，这就产生了在时间平均上不会消失的干涉项，而现在这种增强和减弱是不稳定的，能量在空间上就平均分配了，总的光强也就自然是两束光的光强之和了。

此外需要注意的是，普通的光源发出的光波是包含不同波长的复色光，在实践中我们可以通过各种方式将光的频率波长在一个小的范围内，形成适合干涉的准单色光。另外一种更为先进的作法是使用激光，它的原理是通过外加辐射控制电子跃迁，从而产生单色性很好的光源。

总而言之，光的干涉的条件相较于机械波或低频的无线电波的干涉要苛刻的多，光的单色性还姑且较容易解决，关键是在于是如何得到了两列稳定的相位差的光波，为此科学家设计了各类干涉装置，其共同思想是，先将同一处光源发出的光波分为两束，然后使两束光经过不同的光路再在空间中重新相遇产生干涉。取决于具体的分光方式，有两种不同的干涉方式
\begin{itemize}
    \item 第一种方法称为\uwave{分波面干涉法}，以\uwave{双缝干涉}为代表。
    \item 第二种方法称为\uwave{分振幅干涉法}，以\uwave{薄膜干涉}为代表。
\end{itemize}
在接下来的两节中，我们将分别介绍这两种干涉方法。

\input{chapter/Chapter13C.tex}
\section{分振幅干涉}
\label{sec:分振幅干涉}
在生活中我们常常会见到这样的现象，在阳光照耀下，肥皂泡或水面上的油膜呈现出复杂的彩色条纹，这很明显就是白光的干涉现象，然而这种情况下，我们并没有看到任何的双缝，这就表明还存在双缝之外可以产生相干光的方法，实际上，这些现象中的干涉来自薄膜的结构。

在开始解释薄膜为何可以产生干涉现象前，我们有必要先补充两项基础知识。

\subsection{光的反射和折射定律}
\label{subsec:光的反射和折射定律}
光的反射定律和折射定律是几何光学的两项基本定律，我们在此将其总结如下。
\begin{OldBoxLaw}[光的反射定律]
    光在介质表面发生反射时，入射角等于反射角
    \begin{Equation}&[]
        \theta_1=\theta_2
    \end{Equation}
\end{OldBoxLaw}
\begin{OldBoxLaw}[光的折射定律]
    光在介质表面发生折射时，入射角与反射角的正弦之比，等同于两者折射率的反比
    \begin{Equation}&[]
        n_1\sin\theta_1=n_2\sin\theta_2
    \end{Equation}
\end{OldBoxLaw}
在习惯上，将折射率较小的介质称为\uwave{光疏介质}，将折射率较大的介质称为\uwave{光密介质}，这一点我们很容易记忆，因为任何介质的折射率都大于$1$，而真空的折射率恰为$1$，显然真空是各种介质中最“稀疏”的了。基于这样的定义，我们可以将\xref{law:光的折射定律}中的数学表达式定性表述为
\begin{itemize}
    \item 如\xref{fig:光疏至光密的折射}，当由光疏介质进入光密介质（例如由空气进入玻璃），即当$n_2>n_1$折射率增大时，折射角会比入射角更小，即$\theta_2<\theta_1$，折射光会向内偏转，方向变得更为垂直。
    \item 如\xref{fig:光疏至光密的折射}，当由光密介质进入光疏介质（例如由玻璃进入空气），即当$n_2<n_1$折射率减小时，折射角会比入射角更大，即$\theta_2>\theta_1$，折射光会向外偏转，方向变得更为水平。
\end{itemize}

\begin{OldFigure}[光的反射定律与折射定律]
    \begin{OldFigureSub}[光的反射]
        \inputfigure[scale=1.2]{Chapter13D_Figure01}
    \end{OldFigureSub}
    \hspace{0.5cm}
    \begin{OldFigureSub}[光疏至光密的折射]
        \inputfigure[scale=1.2]{Chapter13D_Figure02}
    \end{OldFigureSub}
    \hspace{0.5cm}
    \begin{OldFigureSub}[光密至光疏的折射]
        \inputfigure[scale=1.2]{Chapter13D_Figure03}
    \end{OldFigureSub}
\end{OldFigure}

\subsection{光的反射与半波损失}
\label{subsec:光的反射与半波损失}
我们知道，无论是从光疏介质进入光密介质，还是从光密介质进入光疏介质，光在界面上都会同时发生反射和折射，在几何光学中，两类情况发生的反射是完全相同的，反射定律总是适用的，反射角总是等于入射角度。然而，在波动光学中，由于我们将光视为一种波，那么我们在研究光的行为就要考虑其作为波的性质，这就会涉及到反射时可能发生的\uwave{半波损失}的问题
\begin{itemize}
    \item 如\xref{fig:光疏至光密的折射}，当由光疏介质进入光密介质，将发生半波损失（用虚线标识）。
    \item 如\xref{fig:光密至光疏的折射}，当由光密介质进入光疏介质，不发生半波损失（用实线标识）。
\end{itemize}
所谓半波损失就是指反射时波发生了$\pi$的相位跃变，如果对应到波长，就相当于是光额外走过了其半波长$\lambda/2$的空间距离。应注意，只有反射会出现半波损失的问题，而折射不会。

\subsection{薄膜干涉}
\label{subsec:薄膜干涉}
现在让我们来正式的讨论薄膜干涉的原理，如\xref{fig:薄膜干涉}，假设光线以入射角$i$由折射率为$n_1$的介质射向折射率为$n_2$的薄膜，此时光线将会在上表面同时发生反射和折射，这将产生两条光路
\begin{itemize}
    \item 光线1是由入射光线在薄膜上表面反射形成的。
    \item 光线2是由入射光纤在薄膜上表明折射形成的，其折射进入薄膜后，第一步将在薄膜下表面发生反射，第二步将通过折射再次穿过薄膜上表面，离开薄膜并与$1$形成平行光。
\end{itemize}

\begin{OldFigure}[薄膜干涉]
    \inputfigure[scale=0.9]{Chapter13D_Figure04}
\end{OldFigure}

这里的光线$1$和光线$2$虽然原本是平行光，但是我们可以用透镜使其汇聚在某一点上，而我们看到，此处的两条光线在$A$点分离前是完全属于同一光线，因此两者最初的频率和初相位是等同的，而两者在再次相遇前，具有恒定的光程差，因此，具有恒定的相位差，所以这就可以说明由薄膜上表面反射的光和由薄膜下表面反射的光是一组相干光，由于这两束光的能量均来自最初的入射光，而能量以振幅形式表现于波动，因此薄膜干涉也被称为分振幅干涉。

由于在薄膜干涉中发生了两次反射，取决于是否产生半波损失，有四种情况
\begin{OldTable}[薄膜干涉的界面反射条件]{|c|c|c|c|}
<\mc{2}(|c|){界面反射条件相同}&\mc{2}(|c|){界面反射条件不同}\\>
\xcell<c>{\inputfigure[scale=0.8]{Chapter13D_Figure05}}&
\xcell<c>{\inputfigure[scale=0.8]{Chapter13D_Figure06}}&
\xcell<c>{\inputfigure[scale=0.8]{Chapter13D_Figure07}}&
\xcell<c>{\inputfigure[scale=0.8]{Chapter13D_Figure08}}\\
$\delta'=0$&$\delta'=0$&$\delta'=-(\lambda/2)$&$\delta'=+(\lambda/2)$\\
\end{OldTable}
在\xref{tab:薄膜干涉的界面反射条件}中的$\delta'$称为附加光程差，这是比较简单的。而接下来，我们的目标就是要计算两束反射光线因其光路不同导致的光程差$\delta_0$，从而最终计算出薄膜干涉中的总光程差$\delta=\delta_0+\delta'$。

\begin{OldBoxFormula}[薄膜干涉的光程差]
    薄膜干涉中，在薄膜上表面和薄膜下表面反射的两束光的光程差为
    \begin{Equation}&[]
        \delta=2e\sqrt{n_2^2-n_1^2\sin^2 i}+\delta'
    \end{Equation}
    其中$e$为薄膜厚度，$n_1,n_2$分别为薄膜外和薄膜的折射率，$i$为入射角。
\end{OldBoxFormula}

\begin{Proof}
    如\xref{fig:薄膜干涉}所示，作$CH\perp AH$，我们看到两条光线在经过$CH$后直至汇聚前经历的光程是完全相同的，因此两者的光程差仅来自于在$A$点分离后在抵达$CH$前的光路差异，具体的说
    \begin{itemize}
        \item 光线1在$n_1$中经过$AH$，记该段的光程为$\Delta_1$。
        \item 光线2在$n_2$中经过$AB$和$BC$，记该段的光程为$\Delta_2$。
    \end{itemize}
    记薄膜中的折射角为$\gamma$，由于薄膜的厚度$BN=e$是已知的，注意到以下的几何关系
    \begin{Equation}&[1]
        AN=NC=e\tan\gamma
        \qquad
        AB=BC=\frac{e}{\cos\gamma}
    \end{Equation}
    故光程$\Delta_1$可以表示为
    \begin{Equation}&[2]
        \Delta_1=n_1(AH)=n_1(AC\sin i)=2n_1e\tan\gamma\sin i
    \end{Equation}
    故光程$\Delta_2$可以表示为
    \begin{Equation}&[3]
        \Delta_2=n_2(AB+BC)=2n_2\frac{e}{\cos\gamma}
    \end{Equation}
    记由光路差异（而不包含半波损失）产生的光程差为$\delta_0=\Delta_2-\Delta_1$，因而
    \begin{Equation}&[4]
        \delta_0=2n_2\frac{e}{\cos\gamma}-2n_1e\tan\gamma\sin i
    \end{Equation}
    根据\fancyref{law:光的折射定律}
    \begin{Equation}&[5]
        n_1\sin i=n_2\sin\gamma
    \end{Equation}
    以\xrefpeq{5}代换\xrefpeq{4}中第二项的$n_1\sin i$
    \begin{Equation}&[6]
        \delta_0=2n_2\frac{e}{\cos\gamma}-2n_2e\tan\gamma\sin\gamma
    \end{Equation}
    由于$\tan\gamma=(\sin\gamma/\cos\gamma)$，我们可以将$2n_2(e/\cos\gamma)$提出
    \begin{Equation}&[7]
        \delta_0=2n_2\frac{e}{\cos\gamma}\qty(1-\sin^2\gamma)
    \end{Equation}
    由于$\cos\gamma=\sqrt{1-\sin^2\gamma}$，分母将被约去
    \begin{Equation}&[8]
        \delta_0=2n_2e\sqrt{1-\sin^2\gamma}
    \end{Equation}\nopagebreak
    至此$\delta_0$的形式已经简洁了很多，我们可能在疑惑为什么要保留$\sqrt{1-\sin^2\gamma}$而不是用更简单的$\cos\gamma$表示之，这是因为折射角$\gamma$并非已知条件，我们需要将其转化为入射角$i$的表示。\goodbreak

    为此，在\xrefpeq{8}再次代入\xrefpeq{5}，考虑到$\sin^2\gamma=(n_1/n_2)^2\sin^2 i$
    \begin{Equation}&[9]
        \delta_0=2n_2e\sqrt{1-\qty(\frac{n_1}{n_2})^2\sin^2 i}
    \end{Equation}
    将$n_2$置入根号内
    \begin{Equation}&[10]
        \delta_0=2e\sqrt{n_2^2-n_1^2\sin i}
    \end{Equation}
    由此就得到了$\delta_0$的表达式，考虑$\delta=\delta_0+\delta'$即得\xrefpeq{}。
\end{Proof}

\fancyref{fml:薄膜干涉的光程差}讨论的是薄膜干涉的一般性原理，尚未涉及任何的具体的干涉情形，在其中我们看到，薄膜干涉的光程差$\delta=2e\sqrt{n_2^2-n_1^2\sin^2 i}+\delta'$在折射率$n_1,n_2$被确定后是关于薄膜厚度$e$和入射角$i$的二元函数$\delta(e,i)$。对于实际的薄膜干涉过程，这两者可以同时变化，但是通常来说，在实践中我们更多遇到的都是其中一个变量被固定的情形
\begin{itemize}
    \item 如果入射角度被固定，称为\uwave{等厚干涉}，此时薄膜厚度相同处将具有相同相位。
    \item 如果薄膜厚度被固定，称为\uwave{等倾干涉}，此时入射角度相同处将具有相同相位。
\end{itemize}
相比之下，等倾干涉较为复杂，我们下面仅就等厚干涉的两类情形进行介绍。

\subsection{劈形膜}
\label{subsec:劈形膜}
\uwave{劈形膜}是一个厚度沿一个方向线性增加的薄膜，如\xref{fig:劈形膜的干涉原理}所示，也称为\uwave{劈尖}。我们假设劈形膜的上下表面间的夹角$\theta$很小，以至于，当我们使用垂直向下的平行光束照射时，在上表面反射的光线可以认为是垂直反射，在下表面反射的光线尽管由于折射的缘故偏移了一段距离，但仍然可以认为与上表面反射的光线是重合的，两者在薄膜上表面附近交会，发生干涉现象。

\begin{OldFigure}[劈形膜的干涉原理]
    \inputfigure[scale=0.9]{Chapter13D_Figure11}
\end{OldFigure}

\begin{OldBoxFormula}[劈形膜干涉的光程差]
    劈形膜在高度为$e$的上表面处产生的光程差$\delta$为
    \begin{Equation}
        \delta=2en_2+\delta'
    \end{Equation}
\end{OldBoxFormula}
\begin{Proof}
    劈形膜干涉时光线均为垂直入射，因此属于等厚干涉，且$i=0$，故有
    \begin{Equation}*
        \delta=2e\sqrt{n_2^2-n_1^2\sin^2 i}+\delta'=2en_2+\delta'\qedhere
    \end{Equation}
\end{Proof}

劈形膜的干涉现象仅取决于其自身的折射率，外部环境的折射率只会影响界面反射条件是否相同，进而影响附加光程差$\delta'$是否为零，从而决定劈形膜高度为零的边缘是明纹还是暗纹
\begin{OldFigure}[劈形膜的干涉条纹]
    \begin{OldFigureSub}[界面反射条件相同]
        \inputfigure[scale=0.95]{Chapter13D_Figure10}
    \end{OldFigureSub}
    \hspace{0.5cm}
    \begin{OldFigureSub}[界面反射条件不同]
        \inputfigure[scale=0.95]{Chapter13D_Figure09}
    \end{OldFigureSub}  
\end{OldFigure}
\begin{OldBoxProperty}[劈形膜在反射条件相同时的干涉条纹]
    当界面反射条件相同时，劈形膜的边缘表现为明纹，此时附加光程差$\delta'=0$。

    当光程差为半波长的偶数倍时产生明纹，其分布为
    \begin{Equation}
        e=\frac{1}{2n_2}\qty(\frac{\lambda}{2})(2m)\qquad
        m\in\N
    \end{Equation}
    当光程差为半波长的奇数倍时产生暗纹，其分布为
    \begin{Equation}
        e=\frac{1}{2n_2}\qty(\frac{\lambda}{2})(2m-1)\qquad
        m\in\N^{*}
    \end{Equation}
\end{OldBoxProperty}

\begin{OldBoxProperty}[劈形膜在反射条件不同时的干涉条纹]
    当界面反射条件不同时，劈形膜的边缘表现为暗纹，此时附加光程差$\delta'=(\lambda/2)$。

    当光程差为半波长的奇数倍时产生暗纹，其分布为
    \begin{Equation}
        e=\frac{1}{2n_2}\qty(\frac{\lambda}{2})(2m)\qquad
        m\in\N
    \end{Equation}

    当光程差为半波长的偶数倍时产生明纹，其分布为
    \begin{Equation}
        e=\frac{1}{2n_2}\qty(\frac{\lambda}{2})(2m-1)\qquad
        m\in\N^{*}
    \end{Equation}
\end{OldBoxProperty}

在实际的劈形膜干涉中，界面反射条件不同，即带有$(\lambda/2)$的附加光程差且边缘为零级暗纹的情况是较常见的，因为通常产生劈形膜的方式，就是利用两片撑开一个小角度的玻璃片间的空气所形成的空气劈尖，这时空气劈形膜上下的玻璃对其而言均为高折射率材料，因此这属于界面反射条件不同的情况（空气--玻璃、玻璃--空气）。值得注意的是，尽管玻璃片本身好像也可以视作某种厚度均匀的薄膜，但玻璃片实际并不会发生干涉，因为玻璃片很厚以至于从下表面反射的光线额外经历的光程差已经超过单个电子跃迁释放的光波的总长度了，与其在上表面相遇的已经不是原来的光波了。由此可见，薄膜干涉的干涉原理只适用于“薄”膜。

在实际应用中，劈形干涉可以用于检测工件表面的平整度，因为，\empx{劈形膜的干涉条纹实际反映的是其上下表面间的等高线}，只不过其是用明纹和暗纹交替标注罢了，在理想情况下这些明纹和暗纹应该是一系列平行的直线，但如果下表面有凹凸不平之处，那么条纹也将随之弯曲。

\begin{OldBoxProperty}[劈形膜的条纹间距]
    劈形膜中，任意两个相邻的明纹或暗纹间的高度差均为
    \begin{Equation}
        \delt{e}=\frac{\lambda}{2n_2}
    \end{Equation}
    劈形膜上的条纹间距可以表示为（基于$\theta$很小近似）
    \begin{Equation}
        L=\frac{\lambda}{2n_2\sin\theta}=\frac{\lambda}{2n_2\theta}
    \end{Equation}
\end{OldBoxProperty}

通过显微镜观察时，我们能观测的是条纹间距$L$，根据\xref{ppt:劈形膜的条纹间距}
\begin{itemize}
    \item 劈形膜的角度越小，折射率越小，光的波长越长（使用红光），产生的干涉条纹就越稀疏。
    \item 劈形膜的角度越大，折射率越大，光的波长越短（使用蓝光），产生的干涉条纹就越密集。
\end{itemize}

\subsection{牛顿环}
\label{subsec:牛顿环}
牛顿环是另外一类常见的薄膜干涉条纹，它的干涉装置由一个置于平板玻璃上的平凸透镜组成，两者间的空气薄膜将对垂直入射的平行光形成一组等厚干涉条纹，而在这种情况下，在相同半径处具有相同的高度，因此这些条纹是以接触点为圆心的一系列同心圆环，称为\uwave{牛顿环}。

牛顿环中的空气薄膜也可以视作上一小节的劈尖薄膜处理，适用\fancyref{fml:劈形膜干涉的光程差}，只不过此时其上表面由平面变为了球面，因此表示牛顿环的光程差表达式不宜再使用薄膜的高度$e$作为自变量，对于薄膜上的点，应该通过其相对于接触点的半径$r$标识位置。
\begin{OldBoxFormula}[牛顿环的光程差]*
    牛顿环在以接触点为圆心的半径$r$处产生的光程差$\delta$为
    \begin{Equation}
        \delta=\frac{r^2}{R}n_2+\frac{\lambda}{2}
    \end{Equation}
    其中，$R$为透镜曲率半径，$n_2$为空气薄膜的折射率，$\lambda$为光波长。
\end{OldBoxFormula}

\begin{Proof}
    牛顿环的干涉装置如\xref{fig:牛顿环的仪器图}所示
    \begin{OldFigure}[牛顿环的仪器图]
        \inputfigure[scale=0.75]{Chapter13D_Figure13}
    \end{OldFigure}
    牛顿环中光线同样是垂直入射的，属于等厚干涉，适用\fancyref{fml:劈形膜干涉的光程差}
    \begin{Equation}&[1]
        \delta=2en_2+\delta'
    \end{Equation}
    我们现在要做的只不过是将$e$重新用$r$表示，在\xref{fig:牛顿环的仪器图}的$\triangle OHP$中关注到
    \begin{Equation}&[2]
        r^2=R^2-(R-e)^2=2Re-e^2
    \end{Equation}
    我们假定透镜的曲率半径远大于薄膜的高度，即总有$R\gg e$，则
    \begin{Equation}&[3]
        r^2=2Re
    \end{Equation}
    此处的近似实际是等价于将球面透镜视为一个抛物面的透镜，这样
    \begin{Equation}&[4]
        e=\frac{r^2}{2R}
    \end{Equation}
    将\xrefpeq{4}代入\xrefpeq{1}即得
    \begin{Equation}&[5]
        \delta=2\frac{r^2}{2R}n_2+\delta'=\frac{r^2}{R}n_2+\delta'
    \end{Equation}
    在牛顿环中$\delta'=(\pi/2)$，因为其装置的结构属于界面反射条件不同的情形。
\end{Proof}
牛顿环的干涉装置中，我们将空气薄膜的上表面由球面近似为抛物面（即便不近似对此处讨论也一样），因此很容易想象，随着半径的增加，空气薄膜的高度将平方增加，而根据先前的讨论我们知道，等厚干涉的干涉条纹反映的其实是薄膜的等高线，等高线在高度变化快的地方会显得更为密集，故可以预期，\empx{牛顿环的干涉条纹具有内疏外密的特征}，另外由于牛顿环的干涉中带有半波损失，因此，\empx{牛顿环的中央将呈现暗斑}，如\xref{fig:牛顿环的干涉图样}所示。定量分析也是容易的。

\begin{OldBoxProperty}[牛顿环的条纹间距]
    牛顿环表现为一系列内疏外密的同心圆环条纹，且中央为暗斑。

    当光程差为半波长的奇数倍是产生暗环，其分布为
    \begin{Equation}&[1]
        r=\sqrt{\frac{\lambda}{2}\frac{R}{n_2}(2m)}=\sqrt{mR\lambda n_2}\qquad m\in\N
    \end{Equation}
    当光程差为半波长的偶数倍是产生明环，其分布为
    \begin{Equation}&[2]
        r=\sqrt{\frac{\lambda}{2}\frac{R}{n_2}(2m-1)}\qquad m\in\N^{*}
    \end{Equation}
\end{OldBoxProperty}

\begin{OldFigure}[牛顿环的干涉图样]
    \inputfigure[scale=0.75]{Chapter13D_Figure12}
\end{OldFigure}

牛顿环的一个常见的用途就是用于测定透镜的曲率半径$R$，其利用的是\xrefpeq{1}。
\begin{OldBoxFormula}[牛顿环测量透镜曲率半径]
    若在牛顿环中观察到第$m$级暗环的直径为$d_m=2r_m$，那么透镜半径为\footnote[2]{测定直径而不是半径的原因是，前者在显微镜下是更容易实践的，因为准确找到圆心并不容易。}
    \begin{Equation}&[1]
        R=\frac{d_m^2}{4m\lambda n_2}
    \end{Equation}
    若分别测得第$x$级和$x+m$级的暗环的直径分别为$d_x$和$d_{x+m}$，那么\footnote[3]{由于牛顿环的中央为暗斑，有时确定暗环的绝对级数是困难的，而该公式只需要确定两个暗环的相对级数差即可。}
    \begin{Equation}&[2]
        R=\frac{d_{x+m}^2-d_x^2}{4m\lambda n_2}
    \end{Equation}
\end{OldBoxFormula}

\begin{Proof}
    证明是容易的，对于\xrefpeq{1}，应用\xrefpeq[牛顿环的条纹间距]{1}
    \begin{Equation}
        d_m^2=4r_m^2=4mR\lambda n_2
    \end{Equation}
    移项即得
    \begin{Equation}
        R=\frac{d_m^2}{4m\lambda n_2}
    \end{Equation}
    而对于\xrefpeq{2}，效仿以上过程
    \begin{Gather}[6pt]
        d_x^2=4xR\lambda n_2 \xlabelpeq{3}\\
        d_{x+m}^2=4(x+m)R\lambda n_2\xlabelpeq{4}
    \end{Gather}
    将\xrefpeq{4}与\xrefpeq{3}相减，即得
    \begin{Equation}
        d_{x+m}^2-d_x^2=4\qty[(x+m)-x]R\lambda n_2=4mR\lambda n_2
    \end{Equation}
    移项即得
    \begin{Equation}
        R=\frac{d_{x+m}^2-d_x^2}{4m\lambda n_2}
    \end{Equation}
    这就证明了\xrefpeq{1}与\xrefpeq{2}。
\end{Proof}
\input{chapter/Chapter13E.tex}
\section{光栅衍射}
\label{sec:光栅衍射}
广义的说，具有周期性的空间结构或光学性能的衍射屏，统称为\uwave{光栅}。这里我们考虑的光栅的结构如\xref{fig:光栅}所示，其由$N$条等距分布的狭缝构成，区别于单缝衍射，我们称之为多缝衍射。

这里，取透光部分的宽度为$a$，取遮光部分的宽度为$b$，这样光栅的周期就是$d=a+b$，这就是说每经过$d$的宽度光栅的结构就会重复一次，我们将光栅的空间周期$d$定义为\uwave{光栅常量}。

\begin{OldFigure}[光栅]
    \inputfigure[scale=0.9]{Chapter13F_Figure10}
\end{OldFigure}

那么，光栅衍射的结果与单缝衍射的结果有何不同？
\begin{OldTable}[光栅衍射的成因]{|c|c|}
<光缝数量$N=2$&光缝数量$N=4$\\>
\xcell<c>[0.5ex][0.0ex]{\inputfigure[scale=0.6]{Chapter13F_Figure05}}
&
\xcell<c>[0.5ex][0.0ex]{\inputfigure[scale=0.6]{Chapter13F_Figure06}}\\
\hlinelig
\mc{2}(|c|){光栅衍射的结果}\\
\hlinemid
\xcell<c>[0.5ex][0.0ex]{\inputfigure[scale=0.6]{Chapter13F_Figure09}}
&
\xcell<c>[0.5ex][0.0ex]{\inputfigure[scale=0.6]{Chapter13F_Figure09}}\\
\hlinelig
\mc{2}(|c|){衍射部分}\\
\hlinemid
\xcell<c>[0.5ex][0.0ex]{\inputfigure[scale=0.6]{Chapter13F_Figure07}}
&
\xcell<c>[0.5ex][0.0ex]{\inputfigure[scale=0.6]{Chapter13F_Figure08}}\\
\hlinelig
\mc{2}(|c|){干涉部分}\\
\end{OldTable}
在\xref{tab:光栅衍射的成因}中，我们注意到即便是对于最简单的“双缝衍射”，它的结果相较于单缝衍射就已经复杂很多了，更不要说更复杂的“四缝衍射”。为此我们有必要仔细探究光栅衍射的光学原理。

首先，我们揭示一个重要的实验事实，\empx{单缝衍射图案的形状和位置并不会随着单缝的上下移动而变化}，这就是说，即便单缝偏离了原有的中心位置，衍射图案仍然将会保持在原位，衍射的中央光强仍然在原先的中心位置。这个结果可能很难理解，不过，我们可以这样想，我们知道，在夫琅禾费衍射中，凸透镜对于成像的决定性作用。具体的说，对于衍射角为$\phi$的一组衍射光束的成像点，其实完全取决于以角度$\phi$通过凸透镜光心的光线指向何处，这一条光线是被凸透镜的位置所固定的，而上下移动单缝只会影响这一通过光心的光线到底是由单缝上具体哪一点发出的，而这并不会影响任何衍射角$\phi$时的成像位置，因此不会影响衍射图案。

这样，我们就很容易想象，如果每次只令光栅的$N$条狭缝中的某一条单独工作，而将其他狭缝遮蔽，那么它们光屏上独立产生的衍射图样将是完全相同且重合的，与\xref{sec:光的衍射}中我们熟悉的单缝衍射没有任何的差别。然而，如果我们使得光栅的$N$条狭缝同时工作，我们就会简单的得到一个亮了$N$倍的单缝衍射图案吗？显然不会！因为虽然每一个狭缝产生的单缝衍射图样完全重合，但是它们之间同时是相干光源，光程差取决于缝间距，光强叠加时是适用相干叠加原理，只不过此时的干涉是在衍射光强分布的基础上进行的，当$N=2$时的干涉就是我们前面很熟悉的双缝干涉，当$N\neq 2$时我们称为\uwave{多缝干涉}，它的干涉结果会稍微更复杂一些。

这样一来，我们就可以说，\empx{光栅衍射实质是单缝衍射和多缝干涉共同作用的结果}。由此我们就可以基本定性解释光栅衍射的结果了，我们知道，干涉的明纹周期取决于相邻两条狭缝中央的间距即光栅周期$d$，衍射的明纹周期则取决于狭缝宽度$a$，通常来说，狭缝刻的很细，狭缝的缝宽往往是远小于缝间距的，这就指示$d\gg a$，因此，\empx{光栅衍射的光强分布中，干涉是高频部分，衍射是低频部分}，两者会构成类似波动中的拍现象，即最终光栅的光强分布，是低频的衍射光强分布的包络线下，以高频的干涉光强分布的规律快速波动，基于此，定义以下特征。
\begin{itemize}
    \item 以双缝情况为例，当干涉部分达到光强极大值时，我们就称光栅衍射对应位置形成的峰值为\uwave{主极大}，主极大并不仅限于中央的哪一最高的峰值，主极大可能会随着衍射部分的包络线的起伏而上下起伏，但这都无妨，只要复合以上定义，我们就统称为主极大。
    \item 以双缝情况为例，当干涉部分达到光强极小值时，即光强为零时，此时光栅衍射对应位置的光强也必然为零，形成暗纹，\empx{注意，此处暗纹的来源是干涉部分而不是衍射部分}！
    \item 如果讨论的是更一般的多缝情形，那么干涉部分会更复杂一点，但总体而言并没有太大的差别，原先的主极大的位置不会改变，主要影响在于，多缝干涉在两个极大值之间还有一些次级的波动，这就会导致所谓\uwave{次极大}的出现，定量的说，\empx{即任意两个主极大之间有$N-1$个暗纹和$N-2$个次极大}，其中$N$为狭缝数量，因此双缝时没有次极大。
    \item 另外还要一种特别的情况，如果$d$刚好是$a$的整数倍，那么衍射部分的周期和干涉部分的周期是耦合的，由于衍射部分也存在光强为零的位置，这就必然会导致样一种情况的出现，主极大恰好落在了衍射部分光强为零的位置，换言之在该处，主极大的峰值恰好为零。从图像上看，我们就好像是意外的丢失一个主极大一样，因此形象的称之为\uwave{缺级}。
\end{itemize}
以上，我们就对光栅衍射中的光强分布现象有了一个基本的理解，下面，我们先不加证明的列出作为绘图依据的光强分布函数，随后给出一系列光栅衍射的光强分布曲线和仿真图像，而在最后，我们将在我们的知识体系范围内半定量的讨论一下主极大的分布规律及其缺级条件。

\begin{OldBoxFormula}[光栅衍射的光强分布]
    光栅衍射在光屏相对于透镜中心$\phi$处产生的光强为
    \begin{Equation}
        I=I_0\qty(\frac{\sin \alpha}{\alpha})^2\qty(\frac{\sin N\beta}{\sin\beta})^2
    \end{Equation}
    其中，前一部分是\uwave{衍射因子}，后一部分是\uwave{干涉因子}，代换变量$\alpha,\beta$分别满足
    \begin{Equation}
        \alpha=\frac{\pi a}{\lambda}\sin\theta\qquad
        \beta=\frac{\pi d}{\lambda}\sin\theta
    \end{Equation}
    其中，$a$和$d$分别为缝宽度和缝的中央间距（光栅常数），$\lambda$为光的波长。
\end{OldBoxFormula}

在\xref{tab:光栅衍射}中，我们以狭缝数目$N=1,2,3,4,5$绘制光栅衍射的光强分布曲线和仿真图像，并分别考虑$d/a=2$和$d/a=4$两种情形，这将决定在一个包络内光强上下振荡次数的多少。

\DeclareDocumentCommand{\tabline}{ms}
{%
    \xcell<c>[0.2ex][0.0ex]{\inputfigure[scale=0.6]{Chapter13F_Figure02_#1}}
    &
    \xcell<c>[0.2ex][0.0ex]{\inputfigure[scale=0.6]{Chapter13F_Figure01_#1}}\\*
    \xcell<c>[0.0ex][0.0ex]{\inputfigure[scale=0.6]{Chapter13F_Figure04_#1}}
    &
    \xcell<c>[0.0ex][0.0ex]{\inputfigure[scale=0.6]{Chapter13F_Figure03_#1}}\\*
    \hlinelig
    \mc{2}(|c|){$N=#1$}\\
    \hlinemid
}

\begin{OldTableLong}[光栅衍射]{|c|c|}
<光栅常量与光缝宽度之比$d/a=2$&光栅常量与光缝宽度之比$d/a=4$\\>
< >( )
\tabline{1}
\tabline{2}
\tabline{3}
\tabline{4}
\tabline{5}
\end{OldTableLong}

需要注意的是，从理论上可以计算，当光栅的狭缝数为$N$时，在中央的主极大处的光强并不是单个狭缝中央明纹的光强$I_0$，而是大的多的$I_0N^2$，定性上说开放的狭缝越多总的光的能量就相应越多，定量的说我们知道在双缝干涉中，如果两个狭缝的光强均为$I_0$，那么两个狭缝的干涉结果的光强极大值为$4I_0$，此处$N=2$恰有$N^2=4$，这就是$N^2$系数的说明。可以预见的是，当$N\to\infty$时，主极大的光强趋于无穷的同时宽度趋于零，即演化为狄拉克函数。

\begin{OldBoxFormula}[光栅的主极大的分布]
    光栅衍射的主极大的分布如下，其中$\phi$为相对中心的角度
    \begin{Equation}&[]
        \phi=\pm(2m)\frac{\lambda}{2d}\qquad m\in\Z
    \end{Equation}
\end{OldBoxFormula}
\begin{Proof}
    现在证明主极大的公式，所谓主极大，其实就是狭缝间两两的光程差$\delta=d\sin\phi$恰好为半波长的偶数倍，这样各个狭缝在在角度$\phi$下发出的光束就具有完全相同的相位，相干增强
    \begin{Equation}*
        \delta=d\sin\phi=\pm(2m)\frac{\lambda}{2}\qquad m\in\Z
    \end{Equation}
    作$\sin\phi=\phi$的小角度近似，将$d$移项
    \begin{Equation}*
        \phi=\pm(2m)\frac{\lambda}{2d}\qedhere
    \end{Equation}
\end{Proof}

\begin{OldBoxFormula}[光栅主极大的缺级条件]
    光栅主极大的缺级条件为，如果第$m$级主极大满足
    \begin{Equation}&[]
        m=\pm \frac{d}{a}m_0\qquad m_0\in\Z
    \end{Equation}
    那么第$m$级主极大将会构成缺级，其中$a,d$分别为光栅的缝宽度和缝的中央间距。
\end{OldBoxFormula}
\begin{Proof}
    根据\fancyref{fml:光栅的主极大的分布}，当$\phi$满足下式时，干涉部分取到极大值
    \begin{Equation}&[1]
        \phi=\pm(2m)\frac{\lambda}{2d}\qquad m\in\Z
    \end{Equation}
    根据\fancyref{fml:单缝衍射的暗纹分布}，当$\phi$满足下式时，衍射部分取到暗纹
    \begin{Equation}&[2]
        \phi=\pm(2m_0)\frac{\lambda}{2a}\qquad m_0\in\Z^{+}
    \end{Equation}
    假如$\phi$同时满足\xrefpeq{1}和\xrefpeq{2}，那么就有
    \begin{Equation}&[3]
        \pm(2m)\frac{\lambda}{2d}=\pm(2m_0)\frac{\lambda}{2a}
    \end{Equation}
    约去无关的系数（两侧正负号是独立的）
    \begin{Equation}
        \frac{m}{d}=\pm\frac{m_0}{a}
    \end{Equation}
    或者
    \begin{Equation}
        m=\pm\frac{d}{a}m_0
    \end{Equation}
    这就证明了\xrefpeq{}的结论。
\end{Proof}

\input{chapter/Chapter13G.tex}